% !TeX root = ../papers/nc-p3.tex
\section{曲線グラフからの純関係の排除と直線を保つ窓の基底}
\label{sec:curve-graph-pure-exclusion}
\begin{theorem}\label{thm:curve-graph-pure-exclusion}
四辺の射影ファイバーが空または一点で、点グラフが一次元の標準有限Frobenius窓とする。
一辺に純関係があれば全ての整曲線成分は全位置で直線となり、この窓はAS正則でない。
従って直線成分を持つ場合、標準Koszul左右AS正則であることと
全四辺の $\mathbb P R_i$ がそれぞれのSegre多様体と交わらないことは同値である。
また、一次元点グラフの全ての標準正則窓は四テンソルとして半安定である。
\end{theorem}
\begin{proof}
純関係 $r=u\otimes v$ と任意の整曲線成分を選ぶ。
前節の最後の束の議論にはAS性も純関係の一意性も使わない。
その議論によれば曲線は直線となり、反対辺の二直線への制限は
そのグラフの式と $r$ の積となる（第二文字側で恒等零なら語順を逆にする）。
ここで前後に対応する成分は実際の点写像で運ぶ。四移動が成分を置換しても、
得られた二直線は同じ一つの辺のグラフで結ばれているから、
定理\ref{thm:line-component-as-test}の当該反対縮約は階数1である。
同定理の必要性から非正則性が従う。逆に純関係が一つもなく、直線成分が存在すれば、
その軌道の各反対縮約は非零で階数2以下なので常に階数2である。同定理の十分性を使う。

半安定性については定理\ref{thm:all-finite-windows-semistable}の最初の重み評価を使う。
一点と直線の積は辺に含まれないので最低初期重みは $\min(d_i,d_{i+1})$ 以下である。
純関係がないので最高初期重みは $\min(e_i+f_{i+1},f_i+e_{i+1})$ 以下である。
相補辺の完全対合で両者を足し四位置で平均すれば、テンソルの最低重みは
各文字重みの平均の和以下となる。これは純関係のない有限Frobenius候補にも成立する。
\end{proof}

\begin{theorem}\label{thm:four-shift-line-pbw}
標準有限Frobenius窓の各位置に射影直線 $L_i$ を取り、標準四移動で
$L_i$ が $L_{i+4}$ に移るとする。各辺は $L_i$ から $L_{i+1}$ への射影同型の
グラフを含み、その各点で当該辺の左右の評価階数がともに3とする。
直線の外の対応には次元・ファイバーの条件を置かない。
$N_i=L_i^\perp$ を二次元文字空間とする。このとき
\[
 R_i\cap (N_i\otimes V_{i+1})=R_i\cap(V_i\otimes N_{i+1})
 =k C_i\subset N_i\otimes N_{i+1}.
\]
窓が標準Koszul左右AS正則となるための必要十分条件は四本の $C_i$ が全て階数2であることである。
その場合、全整数位置に整合する基底で、各辺の六関係は四文字の六つの下降対を先頭とする
二次Gröbner基底になる。
\end{theorem}
\begin{proof}
直線上の同型に沿って二切断を同定する。反対の二直線における制限はグラフの行列式の倍で、
全て零なら評価階数3に反する。従って完全対合により制限四テンソルは
非零二文字関係 $C_i$ とこの行列式の積である。両端で隣接する点を代入すれば
$C_i\in N_i\otimes N_{i+1}$ である。ここでは四移動が指定直線を保つことを使った。

$L_i$ の点 $p$ における評価 $6\mathcal O\to V_{i+1}\otimes\mathcal O(1)$
の階数は3で、余核は次点による $\mathcal O(2)$ である。核 $K$ は階数3、次数 $-2$。
前点で $w$ を縮約しさらに位置 $i+2$ を縮約すれば、$4\mathcal O\to K(1)$ が得られる。
像は隣接関係により $K(1)$ に入り、階数は前点での反対端の辺評価の階数3である。
ゆえに $K(1)$ は大域生成であり、$K$ は自明束の部分束なので、
分解型は $(0,-1,-1)$ である。従って $H^0(K)$ は一次元で第一の交叉を与える。逆順にも同じである。

必要性は単純加群の解消から直接確認できる。直線への商を $D$ とする。
各位置で基底を運べば $D$ は二変数多項式環の反復である。
ASの場合、左単純のKoszul解消に右直線加群 $e_jD$ をテンソルしてTorを求める。
定理\ref{thm:line-component-as-test}の切断計算は指定直線上の評価階数 $1,3,3,1$ だけを使う。
ゆえに $e_jD$ の右極小自由解消は線型で個数は $1,2,1$、その二次関係が $C_j$ である。
右単純の $D$ 上の二変数解消を、各項のこれらの $A$ 上の自由解消で持ち上げる。
具体的には末尾の単射を比較写像で持ち上げて錐を取り、その余核から $e_iD$ への単射を
再び持ち上げて錐を取る。全ての第 $h$ 項が内部次数 $i+h$ で生成されるから得られた解消は極小である。
最終写像の四係数のうち二つは $e_{i+2}D$ の解消の最終写像から来る。
$C_{i+2}$ が階数1ならこの二係数は比例し、四係数の包は三次元以下となる。
一方、最終写像の双対の余核はAS性により単純で、その一次成分を消すには四文字の全支持が必要である。
全ての $i$ に適用する。

十分性のため $N_i$ の補空間 $T_i$ を取り、四整数ごとに基底が戻ることは要求せず、
$N_i$ の基底を順に運んで $C_i$ を二文字の交換子とする。
$T_i$ は直線の切断で基底を運ぶ。$R_i$ の $T_i\otimes T_{i+1}$ への射影は接線文字の交換子一本を張る。
この射影の五次元核から $T_i\otimes N_{i+1}$ への射影は核 $kC_i$ を持ち全射である。
従って四つの混合下降対を、法線文字から始まる対で書き換えられる。
最後の接線文字の降順対は昇順対と接線文字の少ない対で書き換えられる。
同じ長さでは接線文字の個数を先に、次に法線二文字、接線二文字の順の辞書順で比較する。
全ての簡約は小さい語に移る。長さ3の昇順語は20個で、
実際の三次元数も $4^3-2\cdot4\cdot6+4=20$ である。
従って全ての三次曖昧さが解消する。二次先頭対の包含曖昧さはなく、
ダイヤモンド補題により全長で昇順語が基底となる。二次PBW性からKoszul性、
所与の標準双対Frobenius性から左右AS性が従う。
\end{proof}
